% Generated draft: requires mathematical and source review.
\documentclass[11pt]{article}
\usepackage[margin=25mm]{geometry}
\usepackage{fontspec}
\setmainfont{Latin Modern Roman}
\usepackage{amsmath,amssymb,mathtools,unicode-math}
\setmathfont{Latin Modern Math}
\usepackage{enumitem,needspace,xurl,hyperref}
\setlength{\parindent}{0pt}
\setlength{\parskip}{5pt}
\setlength{\emergencystretch}{3em}
\newcommand{\N}{\mathbb{N}}
\newcommand{\Z}{\mathbb{Z}}
\newcommand{\Q}{\mathbb{Q}}
\newcommand{\R}{\mathbb{R}}
\newcommand{\C}{\mathbb{C}}
\newcommand{\F}{\mathbb{F}}
\newcommand{\A}{\mathbb{A}}
\newcommand{\PP}{\mathbb{P}}
\newcommand{\E}{\mathbb{E}}
\newcommand{\eps}{\varepsilon}
\newcommand{\dd}{\,\mathrm{d}}
\begin{document}
\section*{UnsolvedMath ID: 800057}
\textbf{Source problem code:} catalogue entry
\subsection{Definitions and mathematical statement}
Let \(M_n\) denote the smooth manifold underlying complex projective space \(\mathbb{CP}^{n}\), equipped with its usual differentiable structure. Here \(n\) is the complex dimension, hence a nonnegative integer; the source does not specify its intended range or whether dimension zero is included.

A complex structure on \(M_n\) means a smooth integrable bundle endomorphism \(J:TM_n\to TM_n\) satisfying \(J^2=-\operatorname{id}\). Integrability means that \(J\) is induced by an atlas of complex coordinate charts with holomorphic transition maps. Write
\[
\mathcal{J}(M_n)=\{J:TM_n\to TM_n\mid J\text{ is a smooth integrable complex structure}\}.
\]
A complex structure \(J\) admits a K\"ahler metric if there exists a Riemannian metric \(g\) on \(M_n\) such that
\[
g(JX,JY)=g(X,Y),\qquad \omega(X,Y)=g(JX,Y),\qquad \mathrm{d}\omega=0
\]
for all smooth tangent vector fields \(X,Y\). Define
\[
\mathcal{J}_{\mathrm{K}}(M_n)=\{J\in\mathcal{J}(M_n)\mid (M_n,J)\text{ admits a K\"ahler metric}\}.
\]
\textbf{Problem.} Determine all complex structures on \(M_n\), and in particular determine all those admitting a K\"ahler metric; that is, determine \(\mathcal{J}(M_n)\) and its subset \(\mathcal{J}_{\mathrm{K}}(M_n)\), in the classification sense intended by the source.

\textbf{Unresolved qualifications.} The source does not specify the equivalence relation for this classification. The usual interpretation would identify \(J\) and \(J'\) when there is a diffeomorphism \(f:M_n\to M_n\) satisfying \(\mathrm{d}f\circ J=J'\circ\mathrm{d}f\), equivalently a biholomorphism between the resulting complex manifolds; this interpretation requires confirmation. The source also says ``See Problem 14'' and labels the K\"ahler subquestion \((28*)\). Neither the referenced problem nor an explanation of that label is supplied.
\subsection{Short English statement}
Classify complex structures on the usual smooth manifold underlying \(\mathbb{CP}^{n}\), including the subcollection that admits a K\"ahler metric. The source leaves the dimension range and classification equivalence unspecified. Its cross-reference to Problem 14 and label \((28*)\) remain unresolved.
\subsection{Sources}
\begin{itemize}
\item Supplied local UnsolvedMath registry, AMR Open Problem Lists: permanent ID \(800057\), exact source problem code \texttt{catalogue entry}, title ``Hirzebruch's 1954 Problem List --- Problem 6.'' The registry labels this record \texttt{open}; that label is supplied curation, not an independently verified mathematical assessment.
\item Kotschick, \emph{Updates on Hirzebruch's 1954 Problem List} (2013), Problem 6. Supplied source URL: \url{https://arxiv.org/abs/1305.4623}. The background records access on \(2026\text{-}10\text{-}16\) and extraction by \texttt{arxiv\_latex}; these are supplied provenance statements, not a claim of independent source access or verification.
\item Attribution qualification: the registry lists Kotschick as proposer and 2013 as the proposed year, while the title identifies the problem with Hirzebruch's 1954 list. Both attributions are retained without resolving their distinct historical roles beyond the supplied evidence.
\end{itemize}
\end{document}
